\documentclass[11pt,a4paper]{amsart}
\usepackage[T1]{fontenc}
\usepackage{iftex}
\ifPDFTeX\usepackage[utf8]{inputenc}\fi
\usepackage[margin=1in]{geometry}
\usepackage{amsmath,amssymb,mathtools,microtype}
\usepackage[hidelinks]{hyperref}
\emergencystretch=2em
\hypersetup{pdftitle={Tuza's inequality at maximum average degree 15/2},
 pdfauthor={},pdfsubject={Edge-disjoint triangles and triangle edge covers},
 pdfkeywords={Tuza conjecture, maximum average degree, triangle packing}}
\newtheorem{theorem}{Theorem}
\newtheorem{lemma}[theorem]{Lemma}
\newtheorem{proposition}[theorem]{Proposition}
\newtheorem*{question}{Open assertion}
\DeclareMathOperator{\mad}{mad}
\newcommand{\Hi}{\mathcal H}
\newcommand{\LL}{\mathcal L}
\newcommand{\ind}{\mathbf 1}
\newcommand{\wt}{\widehat w}
\makeatletter
\renewcommand{\@settitle}{\begin{center}\baselineskip14\p@\relax
  \bfseries\@title\end{center}}
\makeatother
\title[Tuza's inequality at maximum average degree 15/2]{Tuza's inequality at maximum average degree 15/2}
\date{}
\begin{document}
\raggedbottom
\begin{abstract}
Let $\nu$ and $\tau$ be the edge packing and edge covering numbers of
the triangles of a finite simple graph. We prove $\tau\le 2\nu$ when
the maximum average degree is at most $15/2$, using Kandasamy's
maximum-degree-ten theorem at equality. We also prove that the
vertices of degree at least eight in any counterexample contain two
distinct, possibly overlapping, copies of $K_1\vee P_4$. A shorter
argument handles the case in which those vertices induce a triangle-free
graph. We state the finite local certificate domains and give an open
simultaneous-packing condition sufficient for the unrestricted conjecture.
\end{abstract}
\maketitle
\markboth{TUZA'S INEQUALITY AT MAXIMUM AVERAGE DEGREE 15/2}
         {TUZA'S INEQUALITY AT MAXIMUM AVERAGE DEGREE 15/2}

\section{Statements and conventions}
All graphs are finite and simple. A \emph{packing} is a set of
pairwise edge-disjoint triangles. A \emph{cover} is an edge set meeting
every triangle. Write $\nu(G)$ for the largest size of a packing and
$\tau(G)$ for the smallest size of a cover. Tuza's conjecture asserts
$\tau(G)\le 2\nu(G)$ for every $G$ and remains open in general.
Put
\[
 \mad(G)=\max_{\varnothing\ne V(J),\,J\subseteq G}
             \frac{2|E(J)|}{|V(J)|},
 \qquad H(G)=G[\{v:d_G(v)\ge 8\}],
\]
with $\mad(\varnothing)=0$. Degrees defining $H(G)$ are taken in $G$.

\begin{theorem}[Maximum average degree]\label{thm:density}
If $\mad(G)\le15/2$, then $\tau(G)\le2\nu(G)$, assuming
Kandasamy's Theorem~1.1 \cite{kandasamy}. For $\mad(G)<15/2$ the
same conclusion follows without that maximum-degree-ten theorem.
Both statements use the local inputs in Section~\ref{sec:local}.
\end{theorem}

\begin{theorem}[Triangle-free high subgraph]\label{thm:free}
If $H(G)$ is triangle-free, then $\tau(G)\le2\nu(G)$.
\end{theorem}

\begin{theorem}[Two high fans]\label{thm:fans}
If $\tau(G)>2\nu(G)$, then $H(G)$ contains two copies of
$K_1\vee P_4$ with distinct seven-edge sets. The copies may share
vertices and edges and may have the same five-vertex set.
\end{theorem}

A fan is an ordinary subgraph. Four distinct rim vertices $x,b,c,z$
form a path in that order and a fifth vertex $a$ sees all four.
Reversing the path gives the same seven edges. Other edges on these
vertices and triangles using exterior vertices are allowed. All three
theorems concern the packing and covering of \emph{every} triangle of
$G$. Theorem~\ref{thm:fans} implies Theorem~\ref{thm:free}, but
Section~\ref{sec:free} gives a proof using fewer finite inputs.

\begin{figure}[ht]
\centering
\setlength{\unitlength}{1pt}
\begin{picture}(168,62)
\qbezier(12,10)(36,10)(60,10)
\qbezier(60,10)(84,10)(108,10)
\qbezier(108,10)(132,10)(156,10)
\qbezier(84,49)(48,29.5)(12,10)
\qbezier(84,49)(72,29.5)(60,10)
\qbezier(84,49)(96,29.5)(108,10)
\qbezier(84,49)(120,29.5)(156,10)
\put(12,10){\circle*{3}}\put(60,10){\circle*{3}}
\put(108,10){\circle*{3}}\put(156,10){\circle*{3}}\put(84,49){\circle*{3}}
\put(12,2){\makebox(0,0){$x$}}\put(60,2){\makebox(0,0){$b$}}
\put(108,2){\makebox(0,0){$c$}}\put(156,2){\makebox(0,0){$z$}}
\put(84,57){\makebox(0,0){$a$}}
\end{picture}
\caption{The three specified faces are $abx,abc,acz$. The two
repeated edges $ab,ac$ have different roles in the packing in
Section~\ref{sec:fans}, where the two possible head matchings share $bc$.}
\label{fig:fan}
\end{figure}

Puleo proved the unrestricted density bound $\mad(G)<7$
\cite{puleo}. Gupta \cite{gupta} and Kandasamy \cite{kandasamy}
give maximum-degree bounds. Tsujimura's larger density bounds
require restrictions on triangle overlap, including at most four
triangles on an edge \cite{tsujimura}. Our density hypothesis allows
arbitrary maximum degree and edge triangle multiplicity. In an
all-high minimal violation, Puleo's odd-wheel link lemmas already
give at least five fans at every vertex \cite[\S4]{puleo}. The mixed
high- and low-degree case is the issue in Theorem~\ref{thm:fans}.
The triangle graph of a graph has its triangles as vertices, with
adjacency by a common edge. The perfect-triangle-graph criterion
in \cite[Theorem 18]{trianglegraphs} concerns the triangle graph of
the whole graph. We use the triangle graph of $H(G)$ and control
separately the triangles that meet lower-degree vertices.
We make no claim of historical priority.

\section{Minimal violations and local certificates}\label{sec:local}
Choose a violation $M$ minimal under subgraph inclusion. Each proper
subgraph of $M$ satisfies Tuza. We may do this inside any of our
three classes: average-density bounds pass to subgraphs, and a
vertex high in a subgraph is high in the original graph. An edge
of $H(M)$ therefore remains an edge of the original high subgraph.
Write $\Hi=V(H(M))$, call its members \emph{high}, and call a vertex
of degree five or six \emph{low}. A \emph{seven} is a vertex of degree
seven. Put $r(v)=|N(v)\cap\Hi|$. For a high $h$ write $j=r(h)$.
A vertex $h$ \emph{subsumes} its neighbor $v$ if $N[v]\subseteq N[h]$.
The \emph{link} of $v$ is $M[N(v)]$. The \emph{codegree} of an edge
is its number of common neighbors. Call edges of codegree at least
four \emph{wide}. Puleo's reductions imply robustness: every link
component has at least five vertices. Thus a link of order less
than ten is connected. An admissible seven-link must also fail
the following test. If a matching $P$ admits a set of at most
$|P|$ vertices meeting every link edge outside $P$, pack the
root triangles on $P$ and buy the
matched edges and spokes to the chosen vertices. This reduces
the root by Lemma~\ref{lem:pay}. These are the endpoint link
conditions in the pair enumeration.

\begin{lemma}[Local payment]\label{lem:pay}
Let $W\ne\varnothing$ be a set of roots, let $P$ be a packing, and
let $X\subseteq E(M)$. Suppose $|X|\le2|P|$, every triangle at a
root meets $X$, and every edge of $P$ with both ends outside $W$
belongs to $X$. Then $M$ is not a minimal violation.
\end{lemma}
\begin{proof}
Apply minimality to $M-W-X$. Its packing together with $P$ is
edge-disjoint. Its cover together with $X$ covers $M$. These two
extensions give $\tau(M)\le2\nu(M)$. More generally one may delete a
further edge $e\notin X$ if every triangle on $e$ meets $X$ and every
off-root edge of $P$ lies in $X\cup\{e\}$.
\end{proof}

A rooted neighborhood in the checks below includes all neighbors of
each root, all edges of its link, and all specified intersections.
Thus an exterior vertex cannot add a triangle incident to a root
with a new neighbor. Unknown edges between private neighbors of
different roots are included in the pair checks.

\begin{lemma}[Low vertices and seven-links]\label{lem:structure}
In $M$ the minimum degree is at least five, and the lows form an
independent set. Each five has a link missing at most one edge and
has at least three pairwise adjacent high subsumers. A six either has
an octahedral link with six high neighbors (a \emph{thin} six), or its
link lacks at most one edge and it has at least four high subsumers.
A vertex subsuming a low has degree at least nine. For distinct lows
$x,z$, $|N(x)\cup N(z)|\ge7$; equality forces both to be fives
with clique links.

Every seven has $r\ge4$ and at least five wide spokes. If $r=4$,
its four high neighbors induce a square or a clique. Adjacent
sevens cannot have codegree five or six. At codegree four, their
common four neighbors are high and induce a square or a clique.
Each of the two private neighbors on either side sees all four
common neighbors.
\end{lemma}

The low-vertex inputs use Puleo's robustness and subsumption
reductions \cite[\S2]{puleo} and the low-neighborhood union lemma
\cite[Lemma 3.3]{bfg}. For a thin six adjacent to an eight or nine
one uses Kandasamy's local Lemma~3.5 with its connected-neighborhood
hypothesis \cite{kandasamy}. The small-degree proofs also use
Gupta's rooted reductions \cite{gupta}. A codegree-four result with
additional codegree bounds at the tips alone would not justify the
last sentence of the lemma. For four common tips and two private
vertices at each endpoint, a tip is over budget if it sees more
than five of the other seven labels outside the endpoints.
Counting classes by $0,1,2,3,4$ over-budget tips gives
$1{,}144,1{,}134,715,210,83$ respectively. Only three in the
last group survive the uniform payment test. The following are
the extra finite
inputs used here and in the fan proof.

\begin{center}\small
\begin{tabular}{@{}p{.24\linewidth}p{.69\linewidth}@{}}
\hline
Domain & Enumeration and outcome\\ \hline
Single seven &
The $2^{21}$ labeled seven-vertex links give 1,044 isomorphism
classes. Positive rooted certificates among the reduced cases
include 40 cone certificates. The surviving candidates give the
five-wide-spoke bound.\\[3pt]
Two sevens, four tips &
451,235 labeled endpoint pairs give 3,286 marked classes.
Of these, 3,245 admit transferred certificates and 38 further
classes have direct certificates. The 3,283 positive classes each
pass all 16 private-edge completions, for 52,528 checks. One of the
three survivors has a square on the quartet and two have a clique.\\[3pt]
Mixed nine, three high peers &
64 masks in the remaining marked case admit local payments.
The certificates contain 507 packed triangles and account for
3,010 rooted triangles in the local supports.\\[3pt]
High markings and fan counts &
77 surviving unmarked seven-links have 9,856 high/low markings;
159 put every wide spoke at a high neighbor. Five of those have
zero ordinary high fans and none has exactly one, two, or three.
At a tetrahedron-neighboring seven, 194 of 200 remaining labeled
assignments reduce directly.\\[3pt]
Two sevens, three tips &
The path-type pair at the end of the marked-link reduction admits
one fixed payment on all $2^9=512$ private-edge completions.\\ \hline
\end{tabular}
\end{center}
In a positive class a checker tests triangle membership,
edge-disjointness, the covering condition for all rooted
triangles, and payment for the packed edges left in the residual
graph. In classifications using designated exits, it also checks
their payment bounds under the recorded degree hypotheses. A
surviving signature specifies necessary incidences; it need not
be realizable in a violation. The companion verification guide
identifies the implemented replays and inherited inputs. The
bounded census in Section~\ref{sec:census} is independent of all
three proofs.

\section{Counting high neighbors}\label{sec:count}
Let $n=|V(M)|$, $m=|E(M)|$. Begin with surplus $d(v)-7$ at each
vertex. Each five takes $2/3$ from every high subsumer. Each
nonthin six takes $1/4$ from every high subsumer; each thin six
takes $1/6$ from every neighbor. In Puleo's discharging the resulting
surpluses are nonnegative \cite[\S2, Lemmas 2.7--2.8]{puleo}.
If a five has a clique
link, each of its five neighbors subsumes it, hence is high. The five
has surplus $5(2/3)-2=4/3$. Return $4/15$ to each of these neighbors.
Let $s'(v)$ be the resulting surplus. Independently, assign a high
vertex its own unit of weight and divide each other vertex's unit
equally among its high neighbors. Write $w_0$ for these totals.
Then
\begin{equation}\label{eq:sum}
  s'(v)\ge0,\qquad \sum_v s'(v)=2m-7n,\qquad
  \sum_{h\in\Hi}w_0(h)=n .
\end{equation}
If $h$ has degree $d$, let $\ell$ be its number of low neighbors
and $f$ its number of subsumed fives. All nonhigh neighbors have
$r\ge4$ except fives with $r=3$, whose three high neighbors subsume
them. Consequently
\begin{equation}\label{eq:wzero}
 w_0(h)\le1+(d-j)/4+f/12,\qquad w_0(h)\le1+(d-j)/3 .
\end{equation}
If $h$ subsumes a low, then $\ell\le d-6$. At equality at most two
lows are subsumed: their open neighborhoods lie in a set of seven,
where the three-low local reduction applies: if three are subsumed,
their links are cliques, their complementary pairs in the seven-set
are pairwise disjoint, and the seven-set induces $K_7$. Buying its
21 edges pays for the eleven triangles of \cite[Fig.~4]{bfg} at
these three roots. If $j\le2$, a thin
neighbor or a subsumed six would require more high peers. Two paid
fives would share all three subsumers, forcing a union of size at
most seven and clique links, which give four high peers. Thus the
loss at $j\le2$ is at most $2/3$.

\begin{lemma}\label{lem:seven-five}
A seven adjacent to a five has at least five high neighbors.
If a high has exactly two high peers and they are adjacent, each of
its seven-neighbors also has at least five high neighbors.
\end{lemma}
\begin{proof}
Suppose $p$ is a seven adjacent to a five $x$ with $r(p)=4$.
The edge $px$ has three common neighbors: four would make $p$
subsume $x$. Some wide spoke at $p$ ends at a nonhigh $t$.
It cannot end at a five for the same reason. If $t$ is a seven,
Lemma~\ref{lem:structure} makes $x$ private to $p$ and adjacent
to the four common high tips, again impossible. Thus $t$ is a
nonthin six whose four high subsumers comprise all four high
neighbors of $p$. At least three subsumers of $x$ neighbor $p$,
and hence neighbor $t$. Those three and $p$ lie in
$N(x)\cap N(t)$, giving $|N(x)\cup N(t)|\le7$ for lows of
unequal degree. For the second assertion, a seven's high quartet,
if it exists, contains the high in question. A clique gives it
three high peers; in a square its two high peers would be
nonadjacent.
\end{proof}

Consider a \emph{mixed nine}, a nine $h$ subsuming both a five $x$
and a nonthin six $z$. Low independence and the union lemma give
\begin{equation}\label{eq:union}
 N(x)\cup N(z)=\{h\}\cup(N(h)\setminus\{x,z\}),\quad
 |N(x)\cup N(z)|=8,\quad |N(x)\cap N(z)|=3 .
\end{equation}
In particular $x,z$ are the only low neighbors of $h$ and their
payments give $s'(h)\ge 2-2/3-1/4=13/12$. The four high subsumers
of $z$ give $j\ge3$. If $j=3$ the six uses all four available
highs. The $r(x)\ge4$ case is excluded by the three-element
intersection, and the remaining case is exactly the 64-mask
reduction of Section~\ref{sec:local}. In its labeling, $h,a,b$
are three universal vertices of the ten-vertex display. Its other
vertices are $c,x,z,p,q,r,t$; the required edges among them are
$xp,xq,zc,zr,zt,cr,ct$, and exactly $cp,cq,pr,pt,qr,qt$ are
undecided. The links at $h,x,z$ are exhausted. When a private
seven is deleted, buy its $7-d_{\rm display}(v)$ actual exits.
Thus the 925 internal and 56 exit purchases in the 64 records
cover all triangles at the deleted roots even when two private
vertices share an exterior neighbor. It follows that $j\ge4$.
Let $\LL$ be the mixed nines with $j=4$ and put $S_t=N(t)\cap\Hi$.
At $h\in\LL$, $|S_x\cup S_z|=5$. The three possibilities are
\[
\begin{array}{c|ccc}
 (|S_x|,|S_z|)&(3,4)&(3,5)&(4,4)\\ \hline
 |S_x\cap S_z|&2&3&3\\
 \text{upper bound on }w_0(h)&67/30&131/60&11/5 .
\end{array}
\]
For the first two columns the two sevens adjacent to $x$ have
$r\ge5$ and the third seven has $r\ge4$. For the last column,
$x$ and $z$ contribute at most $1/4$ each and the seven at $x$
contributes at most $1/5$. These estimates give the displayed bounds.

\subsection{Two redistributions}
First modify the fractions sent by a five $x$ with $|S_x|=3$.
Mark those of its three subsumers that are in the $(3,4)$ column.
At any mark the subsumed six sees two labels of $S_x$, and both
subsume that six. If all three were marked, the distinct sixes
would partition three labels into pairs: a marked nine cannot
neighbor a second six by \eqref{eq:union}. At one mark subtract
$1/30$ and add $1/60$ to each unmarked share. At two marks subtract
$1/30$ at each and add $1/15$ to the third. These changes preserve
each unit. Write $\wt$ for the resulting weights.

A $(4,4)$ vertex receives nothing. A $(3,5)$ vertex $g$ receives
at most $1/60$. Indeed, if the other two subsumers are marked with
different sixes, their pairs in $S_x$ must both contain $g$.
This gives $g$ two sixes. If they have the same six, the six of $g$
has five high neighbors including all three members of $S_x$.
At least four of the five subsume it, so one of the two marked
vertices would subsume both sixes. Both cases contradict
\eqref{eq:union}. Thus, for $h\in\LL$,
\begin{equation}\label{eq:limit}
\wt(h)\le
 \frac{67}{30}-\frac1{30}
 =\frac{131}{60}+\frac1{60}=\frac{11}{5}.
\end{equation}
Any high can receive at most $1/15$ per adjacent five with $r=3$.

\begin{lemma}[A recipient outside the exceptional set]\label{lem:recipient}
Every $h\in\LL$ has a high neighbor outside $\LL$ that subsumes a low.
\end{lemma}
\begin{proof}
In column $(3,4)$ look at the three labels in $S_x$. If all were
in $\LL$, their unique adjacent sixes would partition the set into
parts of size at least two. The part through $h$ has size two,
leaving one. All three labels subsume $x$. In column $(4,4)$
partition $S_z$ by the unique fives incident to its marked labels.
The part through $h$ has size three; all parts have size at least
two. This leaves one of the four labels, all of which subsume $z$.

In column $(3,5)$ inspect the five labels of $S_z$. Each member
in $\LL$ has a five whose three high neighbors are in $S_z$ and
subsume it. Five labels cannot partition into triples. If two or
more labels are outside $\LL$, one subsumes $z$ because at least
four of the five do. If a unique outsider belongs to none of the
triples, the remaining four would partition into triples. Thus in
this last case it subsumes a five. In every column the chosen
outsider neighbors $h$, since $h$ subsumes the focal low.
\end{proof}

Send $1/30$ from each $h\in\LL$ to one recipient furnished by the
lemma. Write $w$ for the new weight. At a sender,
$w\le11/5-1/30=13/6\le2s'$. A recipient subsumes a low and can
receive at most $j/30$. If $f_3$ counts its adjacent fives with
$r=3$, the first redistribution gives the bound
\begin{equation}\label{eq:sharp}
 \wt\le1+(d-j)/4+3f_3/20-\ind_{\{f_3>0\}}/10 .
\end{equation}
Indeed each such five adds at most $1/3-1/4+1/15=3/20$ above
the baseline. Its two seven-neighbors both see the subsumer, and
each saves at least $1/4-1/5$ by Lemma~\ref{lem:seven-five}.
We count this saving only once. The right side of
\eqref{eq:sharp} plus $j/30$ decreases with $j$ for fixed $d,f_3$.

\subsection{Donor bounds and the endpoint}
For transparency the following table gives the intermediate
ratios \emph{before} the second transfer. The two paired entries
are in order of increasing $j$.
\begin{center}\small
\begin{tabular}{@{}ll@{}}\hline
Degree and condition & Upper bound on $\wt/s'$\\ \hline
Eight & $2$\\
Nine, no subsumed five or two clique-link fives & $5/3$\\
Nine, only one paid five, $j=2$ or $j\ge3$ & $39/20,\ 159/80$\\
Nine, mixed, $j\ge5$ or $j=4$ & $129/65,\ 132/65$\\
Ten, two subsumed fives, $\ell\le3$, $j=3$ or $j\ge4$ & $183/100,\ 168/85$\\
Ten, every other case & $191/105$\\
Every degree at least eleven & $69/35$\\ \hline
\end{tabular}
\end{center}
Here are the details that give the table and allow the additional
$j/30$. An eight has no thin neighbor and subsumes no low.
The degree-eight link reduction gives $j\ge4$ without a global
minimum-degree hypothesis. To see the finite test, mark five
nonhigh slots of its full link $J$, choose a subset $S$ of them,
and put $R=V(J)\setminus S$. Buy the edges within each side of
the partition $\{h\}\cup S,R$ of the cone, and all exits at $S$.
The cost is at most $7|S|+2|E(J[R])|-|E(J)|$.
After earlier cone reductions, the 203 retained links give 11,368
placements. Of these 3,308 already force a mark to have degree
eight; 8,060 pass the displayed bound against twice a cone packing.
Thus $w_0\le1+4/4=2$ at an eight. It sends and receives no new
weight, and $s'=1$. At degree nine,
there are no thin neighbors and at most two paid lows. With no
subsumed five, the respective $(w_0,s')$ bounds for zero, one,
or two payments are $(13/4,2)$, $(5/2,7/4)$ and $(5/2,3/2)$.
For two subsumed fives their neighborhoods are in one eight-set.
If the union has size eight their intersection consists of $h$
and a second vertex $k$. On each side the only possible missing
edge in the five-link is absent, joins $k$ to a private vertex,
or joins two private vertices. Taking the two sides in either
order yields six types. Their payment sizes $(|P|,|X|)$ in
order $CC,CK,CP,KK,KP,PP$ are
$(8,16),(7,14),(7,14),(7,14),(6,11),(7,14)$, with $C$ a
complete link and $K,P$ the two missing-edge types. Transporting
the six certificates covers $7^2$ labeled link choices and all
512 private cross-edge assignments per choice. All triangles
at $h$ and the two fives are in this ten-vertex display.
Thus the union has size seven and both links are cliques. The six high peers
and the two rebates give $w_0\le33/20$ and $s'\ge6/5$.
In these cases $\wt/s'\le5/3$, $s'\ge1$, and even $9/30$ fits.

If there is just one paid five, $s'\ge4/3$. At $j=2$ the high
peers are its fellow subsumers and adjacent. The other six
nonhigh neighbors are sevens, each with $r\ge5$ by
Lemma~\ref{lem:seven-five}. Hence $w_0\le38/15$ and
$\wt/s'\le39/20$. The margin is at least $1/15=2/30$.
At $j\ge3$ one has $w_0\le31/12$ by \eqref{eq:wzero}.
With the second transfer, \eqref{eq:sharp} gives at most $53/20$
if $f_3=1$ and $13/5$ otherwise, both below $8/3$.

A mixed nine outside $\LL$ has $j\ge5$ and $s'\ge13/12$.
Initially $w_0\le25/12$, with receipt at most $1/15$ at the
first transfer. If $r(x)\ge4$, \eqref{eq:sharp} gives
$w\le1+4/4+5/30=13/6$. If $r(x)=3$, the two sevens $p,q$
at the five force $j=5$ and the eight-element union in
\eqref{eq:union} has six highs. When $r(z)\ge5$,
\[
 w\le 1+\tfrac13+\tfrac15+\tfrac25+\tfrac1{15}
               +\tfrac5{30}=\tfrac{13}{6}.
\]
When $r(z)=4$, its two sevens are the same $p,q$. Each sees all
three high subsumers of the five and all four of the six, a union
of six highs. Hence
\[
 w\le 1+\tfrac13+\tfrac14+\tfrac26+\tfrac1{15}
               +\tfrac5{30}=\tfrac{43}{20}<\tfrac{13}{6}.
\]

For reference the elementary degree-ten budget is as follows.
Entries are upper bounds for the original weight and lower bounds
for the rebated surplus, respectively.
\begin{center}\small
\begin{tabular}{@{}lll@{}}\hline
Condition & $w_0$ at most & $s'$ at least\\ \hline
No subsumed low, no thin neighbor & $13/3$ & $3$\\
No subsumed low, some thin neighbor & $3$ & $2$\\
Some subsumed low, $j\le2$ & $11/3$ & $7/3$\\
$j\ge3,\ f\le1,\ \ell\le3$ & $17/6$ & $11/6$\\
$j\ge3,\ f\le1,\ \ell=4$ & $17/6$ & $7/4$\\
$j\ge3,\ f=2,\ \ell\le3$ & $35/12$ & $17/12$\\
$j\ge3,\ f=2,\ \ell=4$ & $35/12$ & $28/15$\\
Three subsumed fives & $7/3$ & $23/15$\\ \hline
\end{tabular}
\end{center}
Indeed a thin neighbor forces $j\ge4$. For a subsumer
$\ell\le4$, and at $\ell=4$ at most two lows are paid as
subsumed lows; the others cost at most $1/6$ each. Two paid fives
there have clique links and yield two rebates. When three fives
are paid, their neighborhoods are five-subsets of the same
eight-set. Two of their complementary triples intersect, forcing
two clique links and six high peers. These facts and
\eqref{eq:wzero} give the table.

A ten with $f=2$, $\ell\le3$ and $j=3$ cannot also pay a nonthin
six. The three neighborhoods would be in one eight-set with four
highs; the six uses the four highs and two sevens. Each five has
union eight with the six, so both fives contain the other two
sevens and three of the four highs. Their intersection would have
order at least four, contradicting the union lemma. There is no
thin neighbor. Thus $s'\ge5/3$ and $\wt/s'\le183/100$, allowing
$3/30$. At $j\ge4$ use $s'\ge17/12$ and \eqref{eq:sharp} at
$j=4$, $f_3\le2$. It gives
$w\le1+6/4+3/10-1/10+4/30=17/6$ when $f_3>0$ and less
when $f_3=0$. In the other rows $w_0/s'\le34/21$, $f\le3$ and
$s'\ge4/3$. Even adding $1/5$ to the original ratio gives
$191/105$. For $j\le7$ the remaining margin is at least
$76/315>7/30$. For $j\ge8$ there are at most two nonhigh
neighbors, giving $w\le1+2/3+2/15+10/30<8/3$.

Put $k=d-7\ge4$ in larger degrees. When $j\le2$, use
$\wt\le1+d/3+1/15$ and $s'\ge k-2/3$. Their ratio is at most
$71/50$; even after $2/30$ there is strict slack. With no
subsumed low there is no receipt and $s'\ge k-d/6$,
$\wt=w_0\le(k+8)/4$ for $j\ge3$. The resulting ratio is at
most $18/13$. Otherwise $j\ge3$, $\ell\le k+1$. The following
lower bounds on the final margin include all possible receipts:
\[
\begin{array}{l|c|c}
 \text{condition}&s'\text{ at least}&2s'-w\text{ at least}\\ \hline
 \ell=k+1 &(5k-7)/6 &(85k-284)/60\\
 \ell\le k,\ f\le k-1,\ f_3>0 &(9k-5f)/12 &(16k-61)/60\\
 \ell\le k,\ f\le k-1,\ f_3=0 &(9k-5f)/12 &(25k-76)/60\\
 \ell=f=k & k/3+8/15 &(16k-23)/60
\end{array}
\]
For the first row $f\le2$ and the loss is at most
$4/3+(k-1)/6$. Use \eqref{eq:sharp} without its saving at
$j=3$. In the next two rows the loss is at most
$2f/3+(k-f)/4$. At $j=3$ the saving $1/10$ cancels $3/30$ if
$f_3>0$. Otherwise use $w\le(k+8)/4+1/10$. The smallest margin
occurs at $f=k-1$. In the last row two complementary triples of
the five-neighborhoods in an eight-set meet, so two clique-link
rebates apply and $j\ge6$. At $j=6$ use
$w\le29/20+2k/5$. Each displayed margin is positive for $k\ge4$.
For the intermediate ratio in the $f\le k-1$ rows, replacing
$w$ by $(k+8)/4+3f/20$ before the second transfer gives
$(24k+111)/(20k+25)\le69/35$; the other rows give smaller
ratios.

\begin{proof}[Proof of Theorem~\ref{thm:density}]
We have $w(h)\le2s'(h)$ for every high $h$, strictly when
$d(h)\ge11$. Both transfers conserve the sum of the weights. By
\eqref{eq:sum}, $n\le2(2m-7n)$, so $2m/n\ge15/2$. This excludes
a violation with $\mad(G)<15/2$. Kandasamy's Theorem~1.1
\cite{kandasamy} gives $\Delta(M)\ge11$ for every violation,
so at least one summand is strict. Now $2m/n>15/2$ and
$\mad(G)\ge2m/n$ proves the inclusive assertion.
\end{proof}

\section{A triangle-free high layer}\label{sec:free}
\begin{proof}[Proof of Theorem~\ref{thm:free}]
In a minimal violation a five or a nonthin six has at least three
pairwise adjacent high subsumers. The high neighbors of a thin six
induce an octahedron. Since each alternative gives a high triangle,
all nonhigh vertices are sevens. If there are none, $M=H(M)$
has no triangles.

Fix a seven $y$. We show that it has a wide edge to another seven.
If all its wide neighbors are high, put $r=r(y)$. In the induced
triangle-free graph $J$ on its $r$ high neighbors, at least five
vertices have degree at least $r-3$: a wide spoke has at least
four tips, of which at most $7-r$ are sevens. Thus $5\le r\le7$.
Adjacent degrees in $J$ sum to at most $r$. At $r=7$ the five
designated vertices would be independent and have degree at most
two. At $r=6$ a vertex of degree at least four would have four
neighbors of degree at most two, leaving too few designated
vertices. Each of the five therefore has degree three and must
see the only seven in the rim. That seven shares five tips with
$y$. At $r=5$, $J$ has minimum degree two. A five-cycle allows
no chord. With a four-cycle the fifth vertex must see exactly two
opposite corners. Thus $J=C_5$ or $K_{2,3}$. In the cycle both
remaining sevens see all five highs. In the bipartite graph both
see the three degree-two highs and at least one sees a fourth.
In both cases $y$ has the required wide seven-neighbor.

The four common neighbors on a wide seven--seven edge must
therefore induce a square $A$, by Lemma~\ref{lem:structure}.
A private high would complete a high triangle on an edge of $A$.
Both endpoints thus have exactly $A$ as high neighborhood and
all their private seven-neighbors see $A$. Applying the same
conclusion at every seven shows that an entire component $B$ of
the sevens has this common high neighborhood. Each vertex of
$B$ has degree three in $B$. A triangle in $B$ would give five
common neighbors to an edge joining sevens, which is excluded.
Consequently $B$ is triangle-free, has at least six vertices
and has $b=|E(B)|\ge9$.

Greedily edge-color $B$ with five colors and assign its four
largest matchings to different vertices of $A$. For each selected
edge pack its triangle with the assigned corner. This gives at
least $4b/5$ edge-disjoint triangles. Buy all edges of $B$ and
all four edges of $A$. These meet every triangle at $B$; every
packed edge surviving deletion of $B$ has been bought.
Since $b+4\le8b/5$, Lemma~\ref{lem:pay} gives a contradiction.
\end{proof}

\section{Two high fans}\label{sec:fans}
Suppose $M$ is a minimal violation with at most one high fan.
The proof also covers a fan-free high layer: in that case omit
the exceptional bases $U,V$ below. In the triangle graph of
$H(M)$ an induced three-vertex path has middle face $abc$ and
outer faces $abx,acz$ on different shared edges. Its ends share
no edge, so $x\ne z$. The union of the faces is exactly a fan,
and its seven edges recover the three faces. This gives a
bijection, even when ambient chords exist.

If an abstract graph has exactly one induced three-vertex path,
that path is an entire component: any vertex adjacent to it
gives a second path through its middle or an end. Hence a unique
fan has no other high triangle sharing any of its seven edges.
Its five vertices induce exactly the pictured fan, as any chord
of the rim makes a further high triangle with $a$. The other
triangle-graph components are cliques. Pairwise edge-intersecting
triangles either have a common edge or lie on a tetrahedron:
after two share $ab$, any third without $ab$ sharing an edge of
both uses their two other vertices and $a$ or $b$. This gives the
fourth face, and no further triangle meets all four. Thus the
other components are singletons, books (all pages on a common
base), or the four faces of a tetrahedron. A high vertex outside
a high tetrahedron $Q$ cannot see two of its corners: it would
form a triangle sharing an edge with the faces of $Q$. Distinct
high tetrahedra share at most one vertex. The unique fan and $Q$
also share at most one vertex, since every adjacent pair of its
five vertices is on a fan face.

\subsection{The remaining seven-links}
The high octahedron, $K_5$ missing at most one edge, and the wheel
over a high square all contain multiple ordinary fans. The first
excludes thin sixes. By Lemma~\ref{lem:structure} a nonthin six
$q$ therefore sees exactly four high vertices, its subsumer
tetrahedron $Q$. Its other neighbors are sevens, nonadjacent
because neither subsumes $q$. Both see all of $Q$.

In the square signature of a wide seven pair, a private high
would form a wheel. A private six cannot have this square among
its four high neighbors, which induce a clique. A private five
would see all four corners and the endpoint seven would subsume
it. Hence the private neighbors are sevens. Each new edge from
an endpoint to a private seven is itself wide with the same
square, so propagation and the cubic payment in
Section~\ref{sec:free} exclude the square case. In a wide
clique-signature pair, no private high sees all of $Q$.
Consequently both endpoints have high neighborhood precisely $Q$
and all their private neighbors see $Q$.

Call a seven \emph{special} if it sees an entire high tetrahedron
and \emph{normal} otherwise. A six cannot neighbor a normal
seven, and its wide spokes cannot end at a seven or a five.
Thus all its wide spokes go to highs. Among the 159 such markings
in Section~\ref{sec:local}, the five with at most one fan all have
zero. Label the complete link by three tips $a,b,c$ and four
ends $u,v,s,t$. It contains the two bases $uv,st$ and all the
tip--end edges, except that $cu$ can be absent. There are no
other link edges. In the complete case all positions are high.
In the deficient case only $c,u$ may be nonhigh, independently.
At a nonhigh slot the common neighbors with the root are three
highs inducing one edge or a three-vertex path.

At a special seven $y$ there can be no further high neighbor
in a wide pair, as seen above. If it has no wide seven partner,
at most four of its required five wide spokes can end in $Q$.
An outside high sees at most one corner of $Q$ and at most the two
other noncorners of the link, so is narrow. A five is also narrow.
Thus a wide spoke ends at a nonthin six $q$. Its high tetrahedron
must be $Q$: it shares at least two corners with $Q$, since
there are only seven slots at $y$. Its other seven-neighbor is
not adjacent to $y$. A second six at $y$ would likewise have
neighborhood $Q,y$ and one further seven. The two sixes would
have neighborhood union of order at most seven, contrary to
Lemma~\ref{lem:structure}. Write $r,s$ for the last two slots.
Neither sees $q$. Each has at most three neighbors in the rim
by narrowness and at least two by minimality. Write $X,Y$ for
their respective neighborhoods on $Q$ and $\epsilon$ for the
edge $rs$. The degree and narrow-spoke conditions give
$2\le |X|+\epsilon,\ |Y|+\epsilon\le3$ with the same two bounds
on each expression. Of the 200 labeled assignments, 194 reduce
and the six exceptions have $\epsilon=1$ and $X,Y$ complementary
pairs of corners.

Each five $x$ must neighbor a seven, since five high neighbors
would induce $K_5$ missing at most one edge. The three high
subsumers of $x$ form a clique among its common neighbors with
that seven. This excludes a normal seven by its marking. At a
special seven on a wide pair, a five would see all four corners,
making that seven its subsumer. In the exceptional $r,s$ slots
it sees only two corners, although all its high subsumers would
have to be there. Thus there are no fives. The remaining slots
$r,s$ in the six exceptional rows are sevens. A normal $r$ would
have three high common neighbors with $y$, whereas $s$ is one
of the three. A special $r$ would have a tetrahedron with two
corners of $Q$, forcing it to be $Q$, of which it misses two
corners. This excludes the exceptional assignments.

Every special seven now lies on a wide clique pair, and its
three nonhigh neighbors see its own $Q$. A component $B$ of
sixes and special sevens consequently attaches to just one
tetrahedron $Q$. Its degrees within $B$ are two and three.
A triangle would contain two sevens with at least five common
neighbors. Thus $B$ is triangle-free and $b=|E(B)|\ge4$.
By Vizing's theorem its edges have a proper four-coloring. Pack
the $b$ colored edges with their assigned corners and one more
face of $Q$. Buy $E(B)\cup E(Q)$, of size $b+6\le2(b+1)$.
Lemma~\ref{lem:pay} removes $B$. Only normal sevens remain.

In the deficient rim an adjacent seven in the endpoint slot $u$
gives a common-tip path; one in the tip slot $c$ gives one
common edge and an isolated common neighbor. For the path type
let the roots be $y,z$, the common-tip path $a-h-b$, their
private bases $uv,st$, and their private tips $c,d$ respectively.
Pack $dtz,chy,byz,uvy,asz$ and buy
\[
 as,\ ay,\ by,\ ch,\ cy,\ dt,\ hz,\ sz,\ tz,\ uv .
\]
The five triangles are edge-disjoint and the ten edges pay for
all 23 rooted triangles and all four surviving packed edges.
The nine optional edges between the private triples belong to
no rooted triangle. This excludes the path type. In the one-edge
type the common edge identifies a whole base at both roots.
At every remaining root the four ends of its two disjoint bases
and two full tips adjacent to all four ends are high. The third
tip may be nonhigh or may miss one end. The only high rim edge
outside the two bases and the sides of their high pages is its
possible single edge to the other end of the deficient base.

\subsection{Cycles, paths and the forest count}
Form a multigraph $D$ with one vertex for each represented base
and one edge for every seven, joining its two vertex-disjoint
bases. It has no loops. No tetrahedron edge is represented,
since its two full tips would force that seven to see the
tetrahedron. Only the repeated edges $U=ab,V=ac$ can be bases
from the fan. There is no edge $UV$ in $D$. Each other base
belongs to its own book component; every side of that book
appears on exactly one high triangle. At $U$ use the matching
$ax,bc$ of two sides, and at $V$ use $az,bc$. These avoid both
bases but share the physical edge $bc$.

Orient a cycle of $k$ edges in $D$ meeting at most one
exceptional base. At each root pack the triangle on its tail
base and two triangles through a matching of full-tip sides of
its head book. If the head is exceptional use its specified
matching. The ordinary heads have distinct books, so the
packed triangles are edge-disjoint, even for a parallel pair.
Buy the $k$ bases, the $2k$ matching edges and all three tip
spokes per root. The entire rim consists of bases and tip--end
edges, so these purchases meet every rooted triangle. They cost
at most $6k$ for $3k$ triangles and buy every off-root packed
edge. Thus every cycle of $D$ must contain both exceptional
bases.

Consider an ordinary book $B$ of degree one or two in $D$.
Choose one incident seven and any one $e$ of its four sides
using the two full tips. Direct its incident path $A\to B$ or
$A\to B\to C$, starting with the chosen root. In the second
case $A\ne C$, as a parallel pair here would reduce. At $B$
choose a matching containing $e$; if needed use a matching at
$C$ with the same exceptional rule. The $k=1$ or $2$ roots
give $3k$ packed triangles. Buy their $k+1$ bases, $3k$ tip
spokes and the $2k-1$ matching sides other than $e$, and delete
$e$ from the residual. The sole high page on $e$ meets the
purchased base $B$. If no seven outside the path sees $e$, the
extended version of Lemma~\ref{lem:pay} works with $6k$
purchases. Consequently each such side requires an outside
seven. At an outside seven the side cannot be a repeated base
or a side of the exceptional fan. If it were a genuine side of
a represented ordinary book, its unique high triangle would
identify its base as $B$, and that seven would already be on
the path. Thus it must be the unique extra edge of the outside
seven's rim. The four different demands at each small ordinary
book, and those at different books, require distinct sevens.

Delete the represented members of $\{U,V\}$ from $D$. The
ordinary vertices form a forest. Each tree $T$ attaches to each
exceptional vertex at most once, since two attachments make a
cycle meeting at most one exception. Let $r\le2$ be their total
number, and let $L,D_2,J$ count vertices of $T$ of degree one,
two and at least three in the full multigraph. Its $k=L+D_2+J$
vertices have degree sum $2(k-1)+r$. Hence $J\le L+r-2\le L$
and $L+D_2\ge1$. The number of root edges accounted for by $T$
satisfies
\[
 k-1+r\le2L+D_2+1\le3(L+D_2).
\]
Every root is counted exactly once since no root joins $U$ to
$V$. Yet the demands require at least $4\sum_T(L+D_2)$ roots.
There can be no seven. In the all-high remainder buy one base
per book and one edge per isolated face while packing one page
from each. On a tetrahedron pack one face and buy two opposite
edges. On a unique fan pack $abx,acz$ and buy $ab,ac$.
Triangle components share no edge. This contradiction proves
Theorem~\ref{thm:fans}.

\section{An open simultaneous-packing condition}\label{sec:four}
Call $L$ \emph{hereditarily admissible and tight} if every
subgraph of $L$ satisfies Tuza and $\tau(L)=2\nu(L)$.
A packing \emph{avoids} a set of edges if it uses none of them.

\begin{question}
Let $L$ be hereditarily admissible and tight. Let $u,a,z,b,y$
be five distinct vertices on a path in that order, with $u,z,y$
pairwise nonadjacent; all other chords are allowed. Put
$F_1=\{ua,az\}$ and $F_2=\{zb,by\}$. Suppose one minimum cover
avoids $F_1\cup F_2$, a maximum packing avoids $F_1$, and a
possibly different maximum packing avoids $F_2$. Is there one
maximum packing avoiding $F_1\cup F_2$?
\end{question}

\begin{proposition}
An affirmative answer for every such $L$ implies Tuza's conjecture.
\end{proposition}
\begin{proof}
In a least violation $G$ with $\nu(G)=p$, deleting an edge $uv$
gives
\[
 2p+1\le\tau(G)\le1+\tau(G-uv)\le1+2\nu(G-uv)\le1+2p .
\]
The deletion is tight of rank $p$. Every edge has a triangle;
for a tip $w$ of $uv$ no maximum packing after deletion can
avoid $uw,wv$ simultaneously, since $uvw$ could then be added.
In any hereditarily admissible tight graph $T$ each single edge
can be avoided by a maximum packing: otherwise buying it back
would give $\tau(T)\le1+2(\nu(T)-1)$. Both invariants add over
the shared-edge components of triangles, hence each component
of a tight graph is tight. A blocked pair of legs must lie in
one component, since different components allow the two separate
avoiding packings to be combined.

Choose among such tight component subgraphs of $G$ with a blocked
induced wedge one of minimum packing number $s$, and call it $O$.
Rename the wedge $u-z-y$. Delete its two legs to form $L$. Every
maximum of $O$ used at least one leg, so $\nu(L)\le s-1$;
buying both legs back gives $2s\le2+\tau(L)\le2+2\nu(L)$.
Thus $L$ is tight of rank $s-1$. Fix a minimum cover of $L$.
On either deleted leg there is a returning triangle with both
old edges unbought; otherwise the other leg alone would extend
the cover of $L$ to one of $O$ of size $2s-1$. If the two
books offered no choice of distinct outside tips, all these
uncovered pages would use the same tip $a$, and buying $za$
would likewise suffice. Thus there are distinct tips $a,b$,
giving the stated path with four old edges outside this one
cover. Each half is individually avoidable by minimality of $s$
and the component argument. A single maximum avoiding both
would admit both returning triangles in $O$ and reach $s+1$.
This proves sufficiency. The simultaneous-avoidance assertion
itself remains open.
\end{proof}

\section{A bounded check on tight triangle components}\label{sec:census}
A triangle component is a nonempty connected component of the
triangle graph, realized on its active edges and their incident
vertices. Every edge of such a support lies in a triangle,
and the triangles are connected by shared edges. Through ten
vertices the only such supports satisfying $\tau=2\nu$ are
$K_4$ and $K_5$. This confirms an earlier order-ten certificate
package dated 31 July 2026. Its accompanying audit checked saved
packing and cover witnesses without regenerating all order-ten
presentations. The classification concerns positive triangle
components of bounded order.

 
The later enumeration uses a catalogue of \emph{all} graphs of
order $n-1$ up to isomorphism, since deleting a vertex from a
triangle component need not preserve triangle connectivity.
It adds one vertex to each parent and considers every neighborhood
of the new vertex. It removes
edges lying in no triangle and retains a spanning single
component. At order nine the catalogue starts from the 12,346
eight-vertex graphs. Equitable refinement followed by permutations
within color classes yields 274,642 canonical records and retains
64 symmetric presentations as a fallback. Repeated presentations
preserve coverage. To certify
$\tau<2\nu$, the program first compares the complement of a
maximum cut, an explicit cover, with twice a greedy packing.
When necessary it computes the exact packing number, followed
by a cover decision if still unresolved. Separate certificates
cover the orders through eight. Regenerating the nine-vertex
catalogue gives 274,706 records. The subsequent unpartitioned
runs gave the following numbers.
\begin{center}\small
\begin{tabular}{@{}rrrrr@{}}\hline
Order & Surviving augmentations & Exact packings & Exact covers & Tight hits\\ \hline
9 & $1{,}111{,}668$ & $6{,}692$ & 0 & 0\\
10 & $63{,}572{,}957$ & $111{,}898$ & 0 & 0\\ \hline
\end{tabular}
\end{center}
Each order-ten parent admits 512 neighborhoods. The second
column counts surviving augmentation instances, with possible
isomorphic or identical supports. The counts prove no
classification at larger orders.

\smallskip
\noindent\begin{minipage}{\linewidth}
\emph{Verification.} From the repository root,
\texttt{python3 -B verification/run.py -{}-max-order 9} selects the
ordinary replay. Selecting \texttt{-{}-max-order 10} includes the
full order-ten census. The verification documentation states the
scope and inputs of each implemented check. Certificate and
arithmetic replays do not mechanize the graph-theoretic deductions
in this paper or certify the external maximum-degree theorem.
\end{minipage}

\renewcommand{\bibliofont}{\footnotesize\raggedright}
\bibliographystyle{amsplain}
\bibliography{references}
\end{document}
